\documentclass[12pt,a4paper]{article}

\usepackage{amsmath}
\usepackage{amsthm}
\usepackage{amssymb}

\usepackage[right=1.5in,top=1in,left=1.5in,bottom=1in]{geometry}

\usepackage{enumitem}

\usepackage[all]{xy}

\usepackage{url}

\let\emptyset\varnothing

%\usepackage{fancyhdr}\usepackage{lastpage}\pagestyle{fancy}
%\renewcommand{\headrulewidth}{0pt}
%\fancyfoot[C]{Page \thepage \hspace{1pt} of \pageref{LastPage}}

\begin{document}

\newpage

\noindent\textsc{\hfill Constructing $\mathbb{Z}$ \hfill} \\
\noindent\textsc{9 November 2023\hfill P. Beall}\vspace{0.5em}
\hrule\vspace{0.5em}
\begin{abstract}
	Integers have applications in some areas of math. The integers form the smallest group containing the natural numbers. I will present a construction of the integers from the naturals, and if time permits, I will briefly outline a construction of the rationals from the integers.
\end{abstract}

Let $\mathbb{N} = \{0,1,2,\dots\}$ be the set of natural numbers. We hopefully know what natural numbers are, and how to add them. The natural numbers have various properties, such as associativity,
\[
	(p+q) + r = p + (q+r)\qquad\text{for all } p,q,r\in\mathbb{N}\,,
\]
and cancellation,
\[
	p + q = p + r \quad\text{implies}\quad q=r\qquad\text{for all } p,q,r\in\mathbb{N}\,.
\]
See e.g. chapter 4 of \cite{Enderton} for a construction of $\mathbb{N}$. Given $\mathbb{N}$, we will now construct $\mathbb{Z}$. Modulo small modifications, this is actually a construction of an abelian group given a commutative monoid.

Define a relation $\sim$ on $\mathbb{N}\times\mathbb{N}$ by
\[
	(m,n)\sim (p,q) \quad\text{ if and only if }\quad m+q = n+p\,.
\]
Suppose $(m,n), (p,q), (s,t)\in\mathbb{N}\times\mathbb{N}$ and $(m,n)\sim(p,q)$ and $(p,q)\sim (s,t)$. Since $m+n=m+n$,
\[
	(m,n)\sim(m,n)
\]
and $\sim$ is reflexive. Since $m+q=n+p$, we have $n+p=m+q$ and thus
\[
	(p,q)\sim(m,n)\,,
\]
so $\sim$ is symmetric. Since $m+q=n+p$ and $p+t=q+s$, we have
\begin{align*}
	(m+q) + (p+t) &= (n+p) + (q+s) \\
	q+p+m+t &= q+p+n+s \\
	m+t &= n+s
\end{align*}
and thus $(m,n)\sim (s,t)$, so $\sim$ is transitive.

It follows that $\sim$ is an equivalence relation, so it induces a partition of $\mathbb{N}\times\mathbb{N}$. For $(m,n)\in\mathbb{N}\times\mathbb{N}$, let $[m,n]$ denote the equivalence class of $(m,n)$. Let
\[
	\mathbb{Z} = {\mathbb{N}\times\mathbb{N}/\sim} = \{[m,n] \mid (m,n)\in\mathbb{N}\times\mathbb{N}\}\,.
\]
Define the binary operation $+$ on $\mathbb{Z}$ by
\[
	[m,n] + [p,q] = [m+p, n+q]
\]
for all $m,n,p,q\in\mathbb{N}$. If $[m,n] = [m',n']$ and $[p,q]=[p',q']$, then $m+n'=n+m'$ and $p+q'=q+p'$, whence
\begin{align*}
	(m+n') + (p+q')&= (n+m') + (q+p') \\
	(m+p) + (n'+q') &= (n+q) + (m'+p')
\end{align*}
and $[m+p,n+q] = [m'+p',n'+q']$. It follows that $+$ is well-defined in $\mathbb{Z}$.

We can verify that $\mathbb{Z}$ with $+$ has the various properties we want it to have. For instance, associativity in $\mathbb{Z}$ follows from associativity in $\mathbb{N}$: if $[m,n],[p,q],[s,t]\in\mathbb{Z}$, then
\begin{align*}
	([m,n] + [p,q]) + [s,t] &= [m+p, n+q] + [s,t] \\
	&= [(m+p)+s, (n+q)+t] \\
	&= [m+(p+s), n+(q+t)] \\
	&= [m,n] + [p+s, q+t] \\
	&= [m,n] + ([p,q] + [s,t])\,.
\end{align*}
 Observe that $[0,0]$ is the additive identity of $\mathbb{Z}$. Unlike $\mathbb{N}$, every element of $\mathbb{Z}$ has an additive inverse. Suppose $[m,n]\in\mathbb{Z}$. We have $[m+n,m+n] = [0,0]$ since $m+n+0 = m+n+0$. Since
\[
	[m,n] + [n,m] = [m+n,m+n] = [0,0]\,,
\]
$[n,m]$ is the additive inverse of $[m,n]$ by definition and commutativity.

We can write $m$ for $[m,0]$, and write $-m$ for $[0,m]$. Indeed, the map $\iota: \mathbb{N}\rightarrow\mathbb{Z}$ given by $m\mapsto [m,0]$ is an injective monoid homomorphism, and $[0,m]$ is the additive inverse of $[m,0]$.

The group $\mathbb{Z}$ with the map $\iota$ has the following universal property. For any group $G$ and monoid homomorphism $f: \mathbb{N}\rightarrow G$, there is a unique group homomorphism $\overline{f}: \mathbb{Z}\rightarrow G$ such that $f = \overline{f}\circ \iota$.
\[\xymatrix{
	\mathbb{Z} \ar@{.}[r]^{\overline{f}} & G \\
	\mathbb{N} \ar[u]^\iota \ar[ur]_f
}\]
Many groups contain an isomorphic copy of $\mathbb{N}$, but this universal property says $\mathbb{Z}$ is the smallest one, since any group containing an isomorphic copy of $\mathbb{N}$ also contains an isomorphic copy of $\mathbb{Z}$.

Similarly, the ring $\mathbb{Q}$ is the smallest commutative ring with $1$ containing an isomorphic copy of $\mathbb{Z}$: for any commutative ring $R$ with identity and injective ring homomorphism $\phi: \mathbb{Z}\rightarrow R$ such that $\phi(q)$ is a unit whenever $q\neq 0$, there is a injective homomorphism $\overline{\phi}: \mathbb{Q}\rightarrow R$ such that $\overline{\phi}\vert_{\mathbb{Z}} = \phi$.
\[\xymatrix{
	\mathbb{Q} \ar@{.}[r]^{\overline{\phi}} & R \\
	\mathbb{Z} \ar[u] \ar[ur]_{\phi}
}\]

Multiplication in $\mathbb{Z}$ can be defined in terms of multiplication in $\mathbb{N}$, so $\mathbb{Z}$ is a ring. Given $\mathbb{Z}$, we can construct $\mathbb{Q}$. Let $\mathbb{Q}$ be the set of equivalence classes induced by the equivalence relation $\approx$ on $\mathbb{Z}\times\mathbb{Z}$ defined by
\[
	(m,n) \approx (p,q) \quad\text{ if and only if }\quad mq = np\,.
\]
Then, we can define the binary operations on $\mathbb{Q}$ in terms of the binary operations on $\mathbb{Z}$, and verify that any desired properties hold.

\clearpage
\pagestyle{plain}

\bibliographystyle{amsplain}
\begin{thebibliography}{1}
	\bibitem{DummitFoote} D.S. Dummit and R.M. Foote. \emph{Abstract Algebra}, 3rd edition.
	
	\bibitem{Enderton}
	H. B. Enderton, \emph{Elements of set theory}. Academic Press, New York, NY, 1977.
\end{thebibliography}

\end{document}